\centerline{\bf \Huge Турбо тренировка V}
\centerline{\bf \Huge Перезагрузка}

\begin{flushright}
        \scriptsize{всё плохо...?}
\end{flushright}

\problem{1}
Приведите пример такого квадратного трехчлена $P(x)$, что при любом $x$ справедливо равенство

$$
P(x)+P(x+1)+\cdots+P(x+10)=x^2
$$


\problem{2}
На окружности даны точки $A$   и $B$.   На окружности берется произвольная точка $M$   и из середины отрезка $BM$   проводится перпендикуляр к прямой $AM$.   Докажите, что этот перпендикуляр проходит через некоторую точку, не зависящую от выбора точки $M.$

\problem{3}
Анджур Аркадьевич, Расул Владимирович и Очир Владимирович сражаются со Змеем Олегычем. Анджур Аркадьевич каждым своим ударом отрубает половину всех голов и еще одну, Расул Владимирович --- треть всех голов и еще две, а Очир Владимирович - четверть всех голов и еще три. Все трое бьют по одному, в том порядке, в котором считают нужным. Если ни один из них не может ударить из-за того, что число голов получится нецелым, то Змей Олегыч съедает преподов. Смогут ли преподы отрубить все головы $20^{20}$-головому Змею?

\problem{4}
Некоторые клетки доски $100\times100$ покрашены в чёрный цвет. Во всех строках и столбцах, где есть чёрные клетки, их количество нечётно. В каждой строке, где есть чёрные клетки, поставим красную фишку в среднюю по счёту чёрную клетку. В каждом столбце, где есть чёрные клетки, поставим синюю фишку в среднюю по счёту чёрную клетку. Оказалось, что все красные фишки стоят в разных столбцах, а синие фишки — в разных строках. Докажите, что найдётся клетка, в которой стоят и синяя, и красная фишки.